\section{Prioridades y plan de construcción}\label{sec:plan}

\subsection{Resumen de prioridades}
\begin{table}[H]
\centering
\caption{Requisitos por prioridad MoSCoW.}
\begin{tabularx}{\textwidth}{@{}p{2.2cm}p{4.2cm}X@{}}
\toprule
\textbf{Prioridad} & \textbf{Significado} & \textbf{Qué incluye} \\
\midrule
\Must & Imprescindible para la demostración. & Usuarios identificados y roles; reportes de daños y solicitudes de ayuda (con vulnerabilidad); registro de donaciones e inventario con trazabilidad; verificación; priorización con TOPSIS y confianza; plan de asignación optimizado que la entidad aprueba; código de un solo uso validado sin conexión; seguimiento del estado; mapa, listas y alertas; generador de datos simulados. \\
\Should & Importante si hay tiempo. & Donaciones dirigidas y alertas de inventario, tablero de necesidades no cubiertas, comparación con asignación voraz, recálculo del plan, carga de datos oficiales, detección de duplicados y anomalías, formulario en papel, reportes sin conexión, incidencias, registro encadenado, bitácora. \\
\Could & Deseable. & Agrupación en perfiles, escenarios ``qué pasaría si'', fotos, exportación, comprobante de donación, estimación de una donación hipotética. \\
\Wont & No en esta versión. & Procesamiento de pagos, conexión automática con sistemas oficiales, satélites, drones, redes sociales, sensores, aprendizaje supervisado, SMS, verificación de identidad contra bases oficiales. \\
\bottomrule
\end{tabularx}
\end{table}

En total el documento define 157 requisitos: 78 funcionales, 27 de datos, 33 del modelo y 19 no funcionales. De ellos, 85 son \Must, 60 son \Should y 12 son \Could.

\subsection{Plan de 12 horas}
El orden va de lo más urgente a lo menos urgente. Si el tiempo no alcanza, se recorta desde el final de la lista, nunca desde el principio.

\begin{table}[H]
\centering
\caption{Plan de construcción propuesto.}
\begin{tabularx}{\textwidth}{@{}clXc@{}}
\toprule
\textbf{Bloque} & \textbf{Horas} & \textbf{Trabajo} & \textbf{Prioridad} \\
\midrule
1 & 1,5 & Entorno, esquema de la base de datos, carga de las proyecciones de viviendas y esqueleto desplegado en Railway. & \Must \\
2 & 1,5 & Generador de datos simulados: barrios, eventos, usuarios ficticios, reportes, solicitudes, donantes, donaciones e inventario. & \Must \\
3 & 2,0 & Autenticación con roles, registro con identificación y consentimiento, y formularios de reporte, de solicitud (con vulnerabilidad) y de donación. & \Must \\
4 & 1,5 & Verificación, recepción de donaciones e inventario, asignación, código de un solo uso con validación local y estados de seguimiento. & \Must \\
5 & 1,5 & TOPSIS para barrios y solicitudes, confianza y alertas, con pesos ajustables. & \Must \\
6 & 1,5 & Optimización: plan de asignación con programación lineal entera, su explicación y la aprobación por la entidad. & \Must \\
7 & 1,5 & Mapa, listas priorizadas, detalle con explicación, inventario y vista de seguimiento. & \Must \\
8 & 1,0 & Accesibilidad, pruebas, ajustes y despliegue final. & \Must \\
\midrule
 & 12,0 & Total & \\
\bottomrule
\end{tabularx}
\end{table}

\textbf{Alcance realista.} Los requisitos \Must{} son demasiados para hacerlos todos de forma completa en 12 horas, y este plan no deja tiempo libre. Se construyen en una versión mínima: un solo flujo de punta a punta (donación, inventario, solicitud, verificación, plan de asignación, aprobación, código, entrega) con datos simulados. Si hay que recortar, se sacan primero los \Should{} (por ejemplo, la carga de datos oficiales, que no está en el plan) y luego los detalles de los \Must{} que no afectan a ese flujo. Las 12 horas son solo de desarrollo: el video de 2 minutos y el reporte de la competencia necesitan tiempo adicional.

\subsection{Estado de construcción}\label{sec:estado}
A 17 de octubre de 2026 el prototipo está desplegado y recorre el ciclo completo con datos simulados. Los ocho bloques del plan se construyeron, y después se agregaron los ajustes que salieron de las primeras pruebas: municipios y barrios reales, plan por paquetes, recorrido, mapa de ruta y control logístico, mapa de necesidades por barrio y la visibilidad de qué pidió cada solicitud.

\begin{table}[H]
\centering
\caption{Estado de los requisitos de mayor peso.}
\begin{tabularx}{\textwidth}{@{}p{4.2cm}X@{}}
\toprule
\textbf{Estado} & \textbf{Requisitos} \\
\midrule
Hecho & Usuarios, roles y consentimiento; reportes y solicitudes (con vulnerabilidad); verificación; donaciones, recepción e inventario por lote; asignación y código de un solo uso validado sin conexión; cadena de auditoría; prioridad por TOPSIS con pesos ajustables y confianza; plan de asignación optimizado con comparación frente a la asignación voraz, equidad, FEFO y donaciones dirigidas; plan por paquetes; mapa, listas, alertas y resumen por municipio; recorrido y panel; mapas de ruta y control logístico; generador de escenarios y escenario de Malambo; reinicio de los datos de demostración. \\
Parcial & AHP (existe el cálculo, sin pantalla propia); accesibilidad (revisión automática hecha; falta una prueba con personas usuarias de lectores de pantalla); entrega sin conexión (probada con el código y sus pruebas; falta una prueba en un Android real); compatibilidad entre tipo de recurso y acceso al barrio (el plan solo avisa que el barrio está aislado). \\
Pendiente & Carga de archivos de datos oficiales (RF-45, RF-46); agrupación en perfiles y detección de anomalías (RF-43, RF-44); escenarios ``qué pasaría si'' (RF-49); formulario impreso y fotos (RF-08, RF-11); exportación y comprobante de donación; análisis de sensibilidad de los pesos (RM-08); posición real del equipo por GPS. \\
\bottomrule
\end{tabularx}
\end{table}

\subsection{Riesgos principales}
\begin{description}[leftmargin=0pt]
  \item[Tiempo.] Una sola persona, 12 horas y un alcance que creció con el seguimiento de entregas, las donaciones y la optimización. Mitigación: construir el flujo mínimo de punta a punta y recortar desde el final.
  \item[Optimización.] El modelo de asignación es lo más difícil de construir y de explicar. Mitigación: empezar con un modelo pequeño y probado, y tener la asignación voraz como respaldo.
  \item[Datos.] Los barrios vienen de OpenStreetMap y varios municipios pequeños solo tienen centros poblados; el reparto de viviendas entre barrios es simulado, y las fuentes oficiales de eventos y daños están por identificar. Mitigación: la suma de los barrios coincide con las viviendas proyectadas por municipio (RD-01), el origen de cada dato se declara (RD-24 y RD-25) y el CSV de barrios se puede ampliar o reemplazar.
  \item[Privacidad.] El sistema maneja datos personales y sensibles. Mitigación: identidades ficticias en el prototipo, acceso por rol y revisión jurídica antes de cualquier uso real.
  \item[Fraude.] No hay una base oficial con la que verificar la identidad, y alguien puede pedir ayuda con una identidad falsa o registrar una donación que no llega. Mitigación: verificación en campo, recepción verificada en el centro de acopio, detección de duplicados y código de un solo uso atado a un equipo.
  \item[Validez de los resultados.] Los datos simulados no prueban el rendimiento real. Mitigación: comunicarlo como validación sintética de la metodología.
\end{description}
